
\documentclass{article}
\usepackage{colortbl}
\usepackage{makecell}
\usepackage{multirow}
\usepackage{supertabular}

\begin{document}

\newcounter{utterance}

\centering \large Interaction Transcript for game `codenames', experiment `risk\_high', episode 1 with nemotron{-}3{-}ultra.
\vspace{24pt}

{ \footnotesize  \setcounter{utterance}{1}
\setlength{\tabcolsep}{0pt}
\begin{supertabular}{c@{$\;$}|p{.15\linewidth}@{}p{.15\linewidth}p{.15\linewidth}p{.15\linewidth}p{.15\linewidth}p{.15\linewidth}}
    \# & $\;$A & \multicolumn{4}{c}{Game Master} & $\;\:$B\\
    \hline

    \theutterance \stepcounter{utterance}  
    & & & \multicolumn{2}{p{0.3\linewidth}}{
        \cellcolor[rgb]{0.9,0.9,0.9}{
            \makecell[{{p{\linewidth}}}]{
                \texttt{\tiny{[GM$|$GM]}}
                {'hidden': {'team': ['rør', 'and', 'jern', 'palme', 'flagermus', 'føniks', 'krone', 'vakuum', 'tjek'], 'innocent': ['model', 'æble', 'center'], 'opponent': ['pool', 'olympus', 'glas', 'server', 'bjørn', 'hest', 'sæl', 'presse'], 'assassin': ['advokat', 'kirke', 'blæksprutte', 'stat', 'bold']}, 'revealed': {'team': {'team': [], 'innocent': [], 'opponent': [], 'assassin': []}, 'opponent': {'team': [], 'innocent': [], 'opponent': [], 'assassin': []}}}
            }
        }
    }
    & & \\ \\

    \theutterance \stepcounter{utterance}  
    & & \multicolumn{4}{p{0.6\linewidth}}{
        \cellcolor[rgb]{0.9,0.9,0.9}{
            \makecell[{{p{\linewidth}}}]{
                \texttt{\tiny{[P1$\langle$GM]}}
                \texttt{Giv en ledetråd med et enkelt ord, der relaterer til nogle af dit holds ord fra listen nedenfor. Du kan vælge, hvor mange og hvilke ord din ledetråd relaterer til.} \\
\texttt{Målet er at få din partner til at gætte alle dit holds ord ud fra de ledetråde, du giver dem, før det andet hold har afsløret alle deres. Hvis din partner gætter på et "modstander"{-} eller et "distraherende" ord, slutter deres gættetur, og modstanderholdet får lov til at tage en tur. Hvis din partner gætter på et "dræberord", slutter spillet, og dit hold taber med det samme.} \\
\texttt{Jo flere ord din ledetråd relaterer til, jo hurtigere kan du vinde spillet, men jo sværere kan associationerne være at gætte for din partner, så vælg dit ledetrådsord med omhu. Ledetrådsord skal være semantisk relateret til målordene, det må ikke være et af ordene på listerne eller indeholde dele af dem.} \\
\\ 
\texttt{Giv altid din ledetråd med et enkelt ord og din kommaseparerede liste over relaterede målord i følgende format, og gør dine svar så korte som muligt, og inkluder aldrig anden tekst end den, der kræves i denne form:} \\
\\ 
\texttt{VEJLEDNING: <ORD>} \\
\texttt{MÅL: <MÅL>} \\
\\ 
\texttt{Dine holdord er: rør, and, jern, palme, flagermus, føniks, krone, vakuum, tjek.} \\
\texttt{Din modstanders holdord er: pool, olympus, glas, server, bjørn, hest, sæl, presse.} \\
\texttt{Distraherende ord er: model, æble, center.} \\
\texttt{Dræberordene er: advokat, kirke, blæksprutte, stat, bold.} \\
            }
        }
    }
    & \\ \\

    \theutterance \stepcounter{utterance}  
    & \multicolumn{4}{p{0.6\linewidth}}{
        \cellcolor[rgb]{0.9,0.9,0.9}{
            \makecell[{{p{\linewidth}}}]{
                \texttt{\tiny{[P1$\rangle$GM]}}
                \texttt{VEJLEDNING: FLYV} \\
\texttt{MÅL: flagermus, føniks} \\
            }
        }
    }
    & & \\ \\

    \theutterance \stepcounter{utterance}  
    & & & \multicolumn{2}{p{0.3\linewidth}}{
        \cellcolor[rgb]{0.9,0.9,0.9}{
            \makecell[{{p{\linewidth}}}]{
                \texttt{\tiny{[GM$|$GM]}}
                \texttt{flyv} \\
            }
        }
    }
    & & \\ \\

    \theutterance \stepcounter{utterance}  
    & & & \multicolumn{2}{p{0.3\linewidth}}{
        \cellcolor[rgb]{0.9,0.9,0.9}{
            \makecell[{{p{\linewidth}}}]{
                \texttt{\tiny{[GM$|$GM]}}
                ['flagermus', 'føniks']
            }
        }
    }
    & & \\ \\

    \theutterance \stepcounter{utterance}  
    & & & \multicolumn{2}{p{0.3\linewidth}}{
        \cellcolor[rgb]{0.9,0.9,0.9}{
            \makecell[{{p{\linewidth}}}]{
                \texttt{\tiny{[GM$|$GM]}}
                {'word': 'flagermus', 'assignment': 'team'}
            }
        }
    }
    & & \\ \\

    \theutterance \stepcounter{utterance}  
    & & & \multicolumn{2}{p{0.3\linewidth}}{
        \cellcolor[rgb]{0.9,0.9,0.9}{
            \makecell[{{p{\linewidth}}}]{
                \texttt{\tiny{[GM$|$GM]}}
                {'word': 'føniks', 'assignment': 'team'}
            }
        }
    }
    & & \\ \\

    \theutterance \stepcounter{utterance}  
    & & \multicolumn{4}{p{0.6\linewidth}}{
        \cellcolor[rgb]{0.9,0.9,0.9}{
            \makecell[{{p{\linewidth}}}]{
                \texttt{\tiny{[GM$\rangle$P2]}}
                \texttt{Giv en kommasepareret liste med op til 2 ord fra følgende liste, som bedst relaterer til eller er tættest forbundet med ordet 'flyv'. Begynd altid din liste af gæt med GÆTS: , og medtag ikke anden tekst i dit svar.} \\
\\ 
\texttt{jern, bjørn, krone, center, æble, model, stat, føniks, vakuum, tjek, palme, pool, flagermus, olympus, presse, kirke, hest, advokat, and, sæl, bold, blæksprutte, rør, server, glas} \\
            }
        }
    }
    & \\ \\

    \theutterance \stepcounter{utterance}  
    & & & \multicolumn{4}{p{0.6\linewidth}}{
        \cellcolor[rgb]{0.9,0.9,0.9}{
            \makecell[{{p{\linewidth}}}]{
                \texttt{\tiny{[GM$\langle$P2]}}
                \texttt{GÆTS: flagermus, and} \\
            }
        }
    }
     \\ \\

    \theutterance \stepcounter{utterance}  
    & & & \multicolumn{2}{p{0.3\linewidth}}{
        \cellcolor[rgb]{0.9,0.9,0.9}{
            \makecell[{{p{\linewidth}}}]{
                \texttt{\tiny{[GM$|$GM]}}
                ['flagermus', 'and']
            }
        }
    }
    & & \\ \\

    \theutterance \stepcounter{utterance}  
    & & & \multicolumn{2}{p{0.3\linewidth}}{
        \cellcolor[rgb]{0.9,0.9,0.9}{
            \makecell[{{p{\linewidth}}}]{
                \texttt{\tiny{[GM$|$GM]}}
                {'word': 'flagermus', 'assignment': 'team'}
            }
        }
    }
    & & \\ \\

    \theutterance \stepcounter{utterance}  
    & & & \multicolumn{2}{p{0.3\linewidth}}{
        \cellcolor[rgb]{0.9,0.9,0.9}{
            \makecell[{{p{\linewidth}}}]{
                \texttt{\tiny{[GM$|$GM]}}
                {'word': 'flagermus', 'assignment': 'team'}
            }
        }
    }
    & & \\ \\

    \theutterance \stepcounter{utterance}  
    & & & \multicolumn{2}{p{0.3\linewidth}}{
        \cellcolor[rgb]{0.9,0.9,0.9}{
            \makecell[{{p{\linewidth}}}]{
                \texttt{\tiny{[GM$|$GM]}}
                {'word': 'and', 'assignment': 'team'}
            }
        }
    }
    & & \\ \\

    \theutterance \stepcounter{utterance}  
    & & & \multicolumn{2}{p{0.3\linewidth}}{
        \cellcolor[rgb]{0.9,0.9,0.9}{
            \makecell[{{p{\linewidth}}}]{
                \texttt{\tiny{[GM$|$GM]}}
                {'word': 'server', 'assignment': 'opponent'}
            }
        }
    }
    & & \\ \\

    \theutterance \stepcounter{utterance}  
    & & & \multicolumn{2}{p{0.3\linewidth}}{
        \cellcolor[rgb]{0.9,0.9,0.9}{
            \makecell[{{p{\linewidth}}}]{
                \texttt{\tiny{[GM$|$GM]}}
                {'hidden': {'team': ['rør', 'jern', 'palme', 'føniks', 'krone', 'vakuum', 'tjek'], 'innocent': ['model', 'æble', 'center'], 'opponent': ['pool', 'olympus', 'glas', 'bjørn', 'hest', 'sæl', 'presse'], 'assassin': ['advokat', 'kirke', 'blæksprutte', 'stat', 'bold']}, 'revealed': {'team': {'team': ['flagermus', 'and'], 'innocent': [], 'opponent': [], 'assassin': []}, 'opponent': {'team': [], 'innocent': [], 'opponent': ['server'], 'assassin': []}}}
            }
        }
    }
    & & \\ \\

    \theutterance \stepcounter{utterance}  
    & & \multicolumn{4}{p{0.6\linewidth}}{
        \cellcolor[rgb]{0.9,0.9,0.9}{
            \makecell[{{p{\linewidth}}}]{
                \texttt{\tiny{[P1$\langle$GM]}}
                \texttt{Ordene flagermus, and blev gættet korrekt. Dit holdkammerats tur endte der.} \\
\texttt{Giv nu endnu en ledetråd, der relaterer til nogle af dine resterende teamord, og en liste over de relaterede målord. Husk at starte din ledetråd med 'VEJLEDNING: ', sæt en ny linje, og start din kommaseparerede liste over målord med 'MÅL: '. Bemærk: Nogle ord er blevet fjernet fra listerne i forhold til tidligere anmodninger.} \\
\\ 
\texttt{Dine resterende holdord er: rør, jern, palme, føniks, krone, vakuum, tjek.} \\
\texttt{Resterende ord for din modstander er: pool, olympus, glas, server, bjørn, hest, sæl, presse.} \\
\texttt{Resterende distraherende ord er: model, æble, center.} \\
\texttt{Resterende dræberord er: advokat, kirke, blæksprutte, stat, bold.} \\
            }
        }
    }
    & \\ \\

    \theutterance \stepcounter{utterance}  
    & \multicolumn{4}{p{0.6\linewidth}}{
        \cellcolor[rgb]{0.9,0.9,0.9}{
            \makecell[{{p{\linewidth}}}]{
                \texttt{\tiny{[P1$\rangle$GM]}}
                \texttt{VEJLEDNING: MØNT} \\
\texttt{MÅL: krone, tjek} \\
            }
        }
    }
    & & \\ \\

    \theutterance \stepcounter{utterance}  
    & & & \multicolumn{2}{p{0.3\linewidth}}{
        \cellcolor[rgb]{0.9,0.9,0.9}{
            \makecell[{{p{\linewidth}}}]{
                \texttt{\tiny{[GM$|$GM]}}
                \texttt{mønt} \\
            }
        }
    }
    & & \\ \\

    \theutterance \stepcounter{utterance}  
    & & & \multicolumn{2}{p{0.3\linewidth}}{
        \cellcolor[rgb]{0.9,0.9,0.9}{
            \makecell[{{p{\linewidth}}}]{
                \texttt{\tiny{[GM$|$GM]}}
                ['krone', 'tjek']
            }
        }
    }
    & & \\ \\

    \theutterance \stepcounter{utterance}  
    & & & \multicolumn{2}{p{0.3\linewidth}}{
        \cellcolor[rgb]{0.9,0.9,0.9}{
            \makecell[{{p{\linewidth}}}]{
                \texttt{\tiny{[GM$|$GM]}}
                {'word': 'krone', 'assignment': 'team'}
            }
        }
    }
    & & \\ \\

    \theutterance \stepcounter{utterance}  
    & & & \multicolumn{2}{p{0.3\linewidth}}{
        \cellcolor[rgb]{0.9,0.9,0.9}{
            \makecell[{{p{\linewidth}}}]{
                \texttt{\tiny{[GM$|$GM]}}
                {'word': 'tjek', 'assignment': 'team'}
            }
        }
    }
    & & \\ \\

    \theutterance \stepcounter{utterance}  
    & & \multicolumn{4}{p{0.6\linewidth}}{
        \cellcolor[rgb]{0.9,0.9,0.9}{
            \makecell[{{p{\linewidth}}}]{
                \texttt{\tiny{[GM$\rangle$P2]}}
                \texttt{Ordene flagermus, and blev gættet korrekt. Din tur endte der.} \\
\texttt{Giv nu endnu en kommasepareret liste med mindst 1 og op til 2 ord fra følgende ordliste, som bedst relaterer til eller er tættest forbundet med ordet 'mønt'. Husk at starte dit svar med GÆTS: . Bemærk: nogle ord er blevet fjernet fra listen sammenlignet med tidligere anmodninger.} \\
\\ 
\texttt{jern, bjørn, krone, center, æble, model, stat, føniks, vakuum, tjek, palme, pool, olympus, presse, kirke, hest, advokat, sæl, bold, blæksprutte, rør, glas} \\
            }
        }
    }
    & \\ \\

    \theutterance \stepcounter{utterance}  
    & & & \multicolumn{4}{p{0.6\linewidth}}{
        \cellcolor[rgb]{0.9,0.9,0.9}{
            \makecell[{{p{\linewidth}}}]{
                \texttt{\tiny{[GM$\langle$P2]}}
                \texttt{GÆTS: krone, presse} \\
            }
        }
    }
     \\ \\

    \theutterance \stepcounter{utterance}  
    & & & \multicolumn{2}{p{0.3\linewidth}}{
        \cellcolor[rgb]{0.9,0.9,0.9}{
            \makecell[{{p{\linewidth}}}]{
                \texttt{\tiny{[GM$|$GM]}}
                ['krone', 'presse']
            }
        }
    }
    & & \\ \\

    \theutterance \stepcounter{utterance}  
    & & & \multicolumn{2}{p{0.3\linewidth}}{
        \cellcolor[rgb]{0.9,0.9,0.9}{
            \makecell[{{p{\linewidth}}}]{
                \texttt{\tiny{[GM$|$GM]}}
                {'word': 'krone', 'assignment': 'team'}
            }
        }
    }
    & & \\ \\

    \theutterance \stepcounter{utterance}  
    & & & \multicolumn{2}{p{0.3\linewidth}}{
        \cellcolor[rgb]{0.9,0.9,0.9}{
            \makecell[{{p{\linewidth}}}]{
                \texttt{\tiny{[GM$|$GM]}}
                {'word': 'krone', 'assignment': 'team'}
            }
        }
    }
    & & \\ \\

    \theutterance \stepcounter{utterance}  
    & & & \multicolumn{2}{p{0.3\linewidth}}{
        \cellcolor[rgb]{0.9,0.9,0.9}{
            \makecell[{{p{\linewidth}}}]{
                \texttt{\tiny{[GM$|$GM]}}
                {'word': 'presse', 'assignment': 'opponent'}
            }
        }
    }
    & & \\ \\

    \theutterance \stepcounter{utterance}  
    & & & \multicolumn{2}{p{0.3\linewidth}}{
        \cellcolor[rgb]{0.9,0.9,0.9}{
            \makecell[{{p{\linewidth}}}]{
                \texttt{\tiny{[GM$|$GM]}}
                \texttt{presse} \\
            }
        }
    }
    & & \\ \\

    \theutterance \stepcounter{utterance}  
    & & & \multicolumn{2}{p{0.3\linewidth}}{
        \cellcolor[rgb]{0.9,0.9,0.9}{
            \makecell[{{p{\linewidth}}}]{
                \texttt{\tiny{[GM$|$GM]}}
                {'word': 'sæl', 'assignment': 'opponent'}
            }
        }
    }
    & & \\ \\

    \theutterance \stepcounter{utterance}  
    & & & \multicolumn{2}{p{0.3\linewidth}}{
        \cellcolor[rgb]{0.9,0.9,0.9}{
            \makecell[{{p{\linewidth}}}]{
                \texttt{\tiny{[GM$|$GM]}}
                {'hidden': {'team': ['rør', 'jern', 'palme', 'føniks', 'vakuum', 'tjek'], 'innocent': ['model', 'æble', 'center'], 'opponent': ['pool', 'olympus', 'glas', 'bjørn', 'hest'], 'assassin': ['advokat', 'kirke', 'blæksprutte', 'stat', 'bold']}, 'revealed': {'team': {'team': ['flagermus', 'and', 'krone'], 'innocent': [], 'opponent': ['presse'], 'assassin': []}, 'opponent': {'team': [], 'innocent': [], 'opponent': ['server', 'sæl'], 'assassin': []}}}
            }
        }
    }
    & & \\ \\

    \theutterance \stepcounter{utterance}  
    & & \multicolumn{4}{p{0.6\linewidth}}{
        \cellcolor[rgb]{0.9,0.9,0.9}{
            \makecell[{{p{\linewidth}}}]{
                \texttt{\tiny{[P1$\langle$GM]}}
                \texttt{Ordet krone blev gættet korrekt. Ordet presse blev gættet, men det er et modstander ord. Dit holdkammerats tur endte der.} \\
\texttt{Giv nu endnu en ledetråd, der relaterer til nogle af dine resterende teamord, og en liste over de relaterede målord. Husk at starte din ledetråd med 'VEJLEDNING: ', sæt en ny linje, og start din kommaseparerede liste over målord med 'MÅL: '. Bemærk: Nogle ord er blevet fjernet fra listerne i forhold til tidligere anmodninger.} \\
\\ 
\texttt{Dine resterende holdord er: rør, jern, palme, føniks, vakuum, tjek.} \\
\texttt{Resterende ord for din modstander er: pool, olympus, glas, bjørn, hest, sæl.} \\
\texttt{Resterende distraherende ord er: model, æble, center.} \\
\texttt{Resterende dræberord er: advokat, kirke, blæksprutte, stat, bold.} \\
            }
        }
    }
    & \\ \\

    \theutterance \stepcounter{utterance}  
    & \multicolumn{4}{p{0.6\linewidth}}{
        \cellcolor[rgb]{0.9,0.9,0.9}{
            \makecell[{{p{\linewidth}}}]{
                \texttt{\tiny{[P1$\rangle$GM]}}
                \texttt{VEJLEDNING: ASKE} \\
\texttt{MÅL: føniks, rør} \\
            }
        }
    }
    & & \\ \\

    \theutterance \stepcounter{utterance}  
    & & & \multicolumn{2}{p{0.3\linewidth}}{
        \cellcolor[rgb]{0.9,0.9,0.9}{
            \makecell[{{p{\linewidth}}}]{
                \texttt{\tiny{[GM$|$GM]}}
                \texttt{aske} \\
            }
        }
    }
    & & \\ \\

    \theutterance \stepcounter{utterance}  
    & & & \multicolumn{2}{p{0.3\linewidth}}{
        \cellcolor[rgb]{0.9,0.9,0.9}{
            \makecell[{{p{\linewidth}}}]{
                \texttt{\tiny{[GM$|$GM]}}
                ['føniks', 'rør']
            }
        }
    }
    & & \\ \\

    \theutterance \stepcounter{utterance}  
    & & & \multicolumn{2}{p{0.3\linewidth}}{
        \cellcolor[rgb]{0.9,0.9,0.9}{
            \makecell[{{p{\linewidth}}}]{
                \texttt{\tiny{[GM$|$GM]}}
                {'word': 'føniks', 'assignment': 'team'}
            }
        }
    }
    & & \\ \\

    \theutterance \stepcounter{utterance}  
    & & & \multicolumn{2}{p{0.3\linewidth}}{
        \cellcolor[rgb]{0.9,0.9,0.9}{
            \makecell[{{p{\linewidth}}}]{
                \texttt{\tiny{[GM$|$GM]}}
                {'word': 'rør', 'assignment': 'team'}
            }
        }
    }
    & & \\ \\

    \theutterance \stepcounter{utterance}  
    & & \multicolumn{4}{p{0.6\linewidth}}{
        \cellcolor[rgb]{0.9,0.9,0.9}{
            \makecell[{{p{\linewidth}}}]{
                \texttt{\tiny{[GM$\rangle$P2]}}
                \texttt{Ordet krone blev gættet korrekt. Ordet presse blev gættet, men det er et modstander ord. Din tur endte der.} \\
\texttt{Giv nu endnu en kommasepareret liste med mindst 1 og op til 2 ord fra følgende ordliste, som bedst relaterer til eller er tættest forbundet med ordet 'aske'. Husk at starte dit svar med GÆTS: . Bemærk: nogle ord er blevet fjernet fra listen sammenlignet med tidligere anmodninger.} \\
\\ 
\texttt{jern, bjørn, center, æble, model, stat, føniks, vakuum, tjek, palme, pool, olympus, kirke, hest, advokat, bold, blæksprutte, rør, glas} \\
            }
        }
    }
    & \\ \\

    \theutterance \stepcounter{utterance}  
    & & & \multicolumn{4}{p{0.6\linewidth}}{
        \cellcolor[rgb]{0.9,0.9,0.9}{
            \makecell[{{p{\linewidth}}}]{
                \texttt{\tiny{[GM$\langle$P2]}}
                \texttt{GÆTS: føniks, vakuum} \\
            }
        }
    }
     \\ \\

    \theutterance \stepcounter{utterance}  
    & & & \multicolumn{2}{p{0.3\linewidth}}{
        \cellcolor[rgb]{0.9,0.9,0.9}{
            \makecell[{{p{\linewidth}}}]{
                \texttt{\tiny{[GM$|$GM]}}
                ['føniks', 'vakuum']
            }
        }
    }
    & & \\ \\

    \theutterance \stepcounter{utterance}  
    & & & \multicolumn{2}{p{0.3\linewidth}}{
        \cellcolor[rgb]{0.9,0.9,0.9}{
            \makecell[{{p{\linewidth}}}]{
                \texttt{\tiny{[GM$|$GM]}}
                {'word': 'føniks', 'assignment': 'team'}
            }
        }
    }
    & & \\ \\

    \theutterance \stepcounter{utterance}  
    & & & \multicolumn{2}{p{0.3\linewidth}}{
        \cellcolor[rgb]{0.9,0.9,0.9}{
            \makecell[{{p{\linewidth}}}]{
                \texttt{\tiny{[GM$|$GM]}}
                {'word': 'føniks', 'assignment': 'team'}
            }
        }
    }
    & & \\ \\

    \theutterance \stepcounter{utterance}  
    & & & \multicolumn{2}{p{0.3\linewidth}}{
        \cellcolor[rgb]{0.9,0.9,0.9}{
            \makecell[{{p{\linewidth}}}]{
                \texttt{\tiny{[GM$|$GM]}}
                {'word': 'vakuum', 'assignment': 'team'}
            }
        }
    }
    & & \\ \\

    \theutterance \stepcounter{utterance}  
    & & & \multicolumn{2}{p{0.3\linewidth}}{
        \cellcolor[rgb]{0.9,0.9,0.9}{
            \makecell[{{p{\linewidth}}}]{
                \texttt{\tiny{[GM$|$GM]}}
                {'word': 'hest', 'assignment': 'opponent'}
            }
        }
    }
    & & \\ \\

    \theutterance \stepcounter{utterance}  
    & & & \multicolumn{2}{p{0.3\linewidth}}{
        \cellcolor[rgb]{0.9,0.9,0.9}{
            \makecell[{{p{\linewidth}}}]{
                \texttt{\tiny{[GM$|$GM]}}
                {'hidden': {'team': ['rør', 'jern', 'palme', 'tjek'], 'innocent': ['model', 'æble', 'center'], 'opponent': ['pool', 'olympus', 'glas', 'bjørn'], 'assassin': ['advokat', 'kirke', 'blæksprutte', 'stat', 'bold']}, 'revealed': {'team': {'team': ['flagermus', 'and', 'krone', 'føniks', 'vakuum'], 'innocent': [], 'opponent': ['presse'], 'assassin': []}, 'opponent': {'team': [], 'innocent': [], 'opponent': ['server', 'sæl', 'hest'], 'assassin': []}}}
            }
        }
    }
    & & \\ \\

    \theutterance \stepcounter{utterance}  
    & & \multicolumn{4}{p{0.6\linewidth}}{
        \cellcolor[rgb]{0.9,0.9,0.9}{
            \makecell[{{p{\linewidth}}}]{
                \texttt{\tiny{[P1$\langle$GM]}}
                \texttt{Ordene føniks, vakuum blev gættet korrekt. Dit holdkammerats tur endte der.} \\
\texttt{Giv nu endnu en ledetråd, der relaterer til nogle af dine resterende teamord, og en liste over de relaterede målord. Husk at starte din ledetråd med 'VEJLEDNING: ', sæt en ny linje, og start din kommaseparerede liste over målord med 'MÅL: '. Bemærk: Nogle ord er blevet fjernet fra listerne i forhold til tidligere anmodninger.} \\
\\ 
\texttt{Dine resterende holdord er: rør, jern, palme, tjek.} \\
\texttt{Resterende ord for din modstander er: pool, olympus, glas, bjørn, hest.} \\
\texttt{Resterende distraherende ord er: model, æble, center.} \\
\texttt{Resterende dræberord er: advokat, kirke, blæksprutte, stat, bold.} \\
            }
        }
    }
    & \\ \\

    \theutterance \stepcounter{utterance}  
    & \multicolumn{4}{p{0.6\linewidth}}{
        \cellcolor[rgb]{0.9,0.9,0.9}{
            \makecell[{{p{\linewidth}}}]{
                \texttt{\tiny{[P1$\rangle$GM]}}
                \texttt{VEJLEDNING: STÅL} \\
\texttt{MÅL: jern, rør} \\
            }
        }
    }
    & & \\ \\

    \theutterance \stepcounter{utterance}  
    & & & \multicolumn{2}{p{0.3\linewidth}}{
        \cellcolor[rgb]{0.9,0.9,0.9}{
            \makecell[{{p{\linewidth}}}]{
                \texttt{\tiny{[GM$|$GM]}}
                \texttt{stål} \\
            }
        }
    }
    & & \\ \\

    \theutterance \stepcounter{utterance}  
    & & & \multicolumn{2}{p{0.3\linewidth}}{
        \cellcolor[rgb]{0.9,0.9,0.9}{
            \makecell[{{p{\linewidth}}}]{
                \texttt{\tiny{[GM$|$GM]}}
                ['jern', 'rør']
            }
        }
    }
    & & \\ \\

    \theutterance \stepcounter{utterance}  
    & & & \multicolumn{2}{p{0.3\linewidth}}{
        \cellcolor[rgb]{0.9,0.9,0.9}{
            \makecell[{{p{\linewidth}}}]{
                \texttt{\tiny{[GM$|$GM]}}
                {'word': 'jern', 'assignment': 'team'}
            }
        }
    }
    & & \\ \\

    \theutterance \stepcounter{utterance}  
    & & & \multicolumn{2}{p{0.3\linewidth}}{
        \cellcolor[rgb]{0.9,0.9,0.9}{
            \makecell[{{p{\linewidth}}}]{
                \texttt{\tiny{[GM$|$GM]}}
                {'word': 'rør', 'assignment': 'team'}
            }
        }
    }
    & & \\ \\

    \theutterance \stepcounter{utterance}  
    & & \multicolumn{4}{p{0.6\linewidth}}{
        \cellcolor[rgb]{0.9,0.9,0.9}{
            \makecell[{{p{\linewidth}}}]{
                \texttt{\tiny{[GM$\rangle$P2]}}
                \texttt{Ordene føniks, vakuum blev gættet korrekt. Din tur endte der.} \\
\texttt{Giv nu endnu en kommasepareret liste med mindst 1 og op til 2 ord fra følgende ordliste, som bedst relaterer til eller er tættest forbundet med ordet 'stål'. Husk at starte dit svar med GÆTS: . Bemærk: nogle ord er blevet fjernet fra listen sammenlignet med tidligere anmodninger.} \\
\\ 
\texttt{jern, bjørn, center, æble, model, stat, tjek, palme, pool, olympus, kirke, advokat, bold, blæksprutte, rør, glas} \\
            }
        }
    }
    & \\ \\

    \theutterance \stepcounter{utterance}  
    & & & \multicolumn{4}{p{0.6\linewidth}}{
        \cellcolor[rgb]{0.9,0.9,0.9}{
            \makecell[{{p{\linewidth}}}]{
                \texttt{\tiny{[GM$\langle$P2]}}
                \texttt{GÆTS: jern, rør} \\
            }
        }
    }
     \\ \\

    \theutterance \stepcounter{utterance}  
    & & & \multicolumn{2}{p{0.3\linewidth}}{
        \cellcolor[rgb]{0.9,0.9,0.9}{
            \makecell[{{p{\linewidth}}}]{
                \texttt{\tiny{[GM$|$GM]}}
                ['jern', 'rør']
            }
        }
    }
    & & \\ \\

    \theutterance \stepcounter{utterance}  
    & & & \multicolumn{2}{p{0.3\linewidth}}{
        \cellcolor[rgb]{0.9,0.9,0.9}{
            \makecell[{{p{\linewidth}}}]{
                \texttt{\tiny{[GM$|$GM]}}
                {'word': 'jern', 'assignment': 'team'}
            }
        }
    }
    & & \\ \\

    \theutterance \stepcounter{utterance}  
    & & & \multicolumn{2}{p{0.3\linewidth}}{
        \cellcolor[rgb]{0.9,0.9,0.9}{
            \makecell[{{p{\linewidth}}}]{
                \texttt{\tiny{[GM$|$GM]}}
                {'word': 'jern', 'assignment': 'team'}
            }
        }
    }
    & & \\ \\

    \theutterance \stepcounter{utterance}  
    & & & \multicolumn{2}{p{0.3\linewidth}}{
        \cellcolor[rgb]{0.9,0.9,0.9}{
            \makecell[{{p{\linewidth}}}]{
                \texttt{\tiny{[GM$|$GM]}}
                {'word': 'rør', 'assignment': 'team'}
            }
        }
    }
    & & \\ \\

    \theutterance \stepcounter{utterance}  
    & & & \multicolumn{2}{p{0.3\linewidth}}{
        \cellcolor[rgb]{0.9,0.9,0.9}{
            \makecell[{{p{\linewidth}}}]{
                \texttt{\tiny{[GM$|$GM]}}
                {'word': 'rør', 'assignment': 'team'}
            }
        }
    }
    & & \\ \\

    \theutterance \stepcounter{utterance}  
    & & & \multicolumn{2}{p{0.3\linewidth}}{
        \cellcolor[rgb]{0.9,0.9,0.9}{
            \makecell[{{p{\linewidth}}}]{
                \texttt{\tiny{[GM$|$GM]}}
                {'word': 'bjørn', 'assignment': 'opponent'}
            }
        }
    }
    & & \\ \\

    \theutterance \stepcounter{utterance}  
    & & & \multicolumn{2}{p{0.3\linewidth}}{
        \cellcolor[rgb]{0.9,0.9,0.9}{
            \makecell[{{p{\linewidth}}}]{
                \texttt{\tiny{[GM$|$GM]}}
                {'hidden': {'team': ['palme', 'tjek'], 'innocent': ['model', 'æble', 'center'], 'opponent': ['pool', 'olympus', 'glas'], 'assassin': ['advokat', 'kirke', 'blæksprutte', 'stat', 'bold']}, 'revealed': {'team': {'team': ['flagermus', 'and', 'krone', 'føniks', 'vakuum', 'jern', 'rør'], 'innocent': [], 'opponent': ['presse'], 'assassin': []}, 'opponent': {'team': [], 'innocent': [], 'opponent': ['server', 'sæl', 'hest', 'bjørn'], 'assassin': []}}}
            }
        }
    }
    & & \\ \\

    \theutterance \stepcounter{utterance}  
    & & \multicolumn{4}{p{0.6\linewidth}}{
        \cellcolor[rgb]{0.9,0.9,0.9}{
            \makecell[{{p{\linewidth}}}]{
                \texttt{\tiny{[P1$\langle$GM]}}
                \texttt{Ordene jern, rør blev gættet korrekt. Dit holdkammerats tur endte der.} \\
\texttt{Giv nu endnu en ledetråd, der relaterer til nogle af dine resterende teamord, og en liste over de relaterede målord. Husk at starte din ledetråd med 'VEJLEDNING: ', sæt en ny linje, og start din kommaseparerede liste over målord med 'MÅL: '. Bemærk: Nogle ord er blevet fjernet fra listerne i forhold til tidligere anmodninger.} \\
\\ 
\texttt{Dine resterende holdord er: palme, tjek.} \\
\texttt{Resterende ord for din modstander er: pool, olympus, glas, bjørn.} \\
\texttt{Resterende distraherende ord er: model, æble, center.} \\
\texttt{Resterende dræberord er: advokat, kirke, blæksprutte, stat, bold.} \\
            }
        }
    }
    & \\ \\

    \theutterance \stepcounter{utterance}  
    & \multicolumn{4}{p{0.6\linewidth}}{
        \cellcolor[rgb]{0.9,0.9,0.9}{
            \makecell[{{p{\linewidth}}}]{
                \texttt{\tiny{[P1$\rangle$GM]}}
                \texttt{VEJLEDNING: FERIE} \\
\texttt{MÅL: palme, tjek} \\
            }
        }
    }
    & & \\ \\

    \theutterance \stepcounter{utterance}  
    & & & \multicolumn{2}{p{0.3\linewidth}}{
        \cellcolor[rgb]{0.9,0.9,0.9}{
            \makecell[{{p{\linewidth}}}]{
                \texttt{\tiny{[GM$|$GM]}}
                \texttt{ferie} \\
            }
        }
    }
    & & \\ \\

    \theutterance \stepcounter{utterance}  
    & & & \multicolumn{2}{p{0.3\linewidth}}{
        \cellcolor[rgb]{0.9,0.9,0.9}{
            \makecell[{{p{\linewidth}}}]{
                \texttt{\tiny{[GM$|$GM]}}
                ['palme', 'tjek']
            }
        }
    }
    & & \\ \\

    \theutterance \stepcounter{utterance}  
    & & & \multicolumn{2}{p{0.3\linewidth}}{
        \cellcolor[rgb]{0.9,0.9,0.9}{
            \makecell[{{p{\linewidth}}}]{
                \texttt{\tiny{[GM$|$GM]}}
                {'word': 'palme', 'assignment': 'team'}
            }
        }
    }
    & & \\ \\

    \theutterance \stepcounter{utterance}  
    & & & \multicolumn{2}{p{0.3\linewidth}}{
        \cellcolor[rgb]{0.9,0.9,0.9}{
            \makecell[{{p{\linewidth}}}]{
                \texttt{\tiny{[GM$|$GM]}}
                {'word': 'tjek', 'assignment': 'team'}
            }
        }
    }
    & & \\ \\

    \theutterance \stepcounter{utterance}  
    & & \multicolumn{4}{p{0.6\linewidth}}{
        \cellcolor[rgb]{0.9,0.9,0.9}{
            \makecell[{{p{\linewidth}}}]{
                \texttt{\tiny{[GM$\rangle$P2]}}
                \texttt{Ordene jern, rør blev gættet korrekt. Din tur endte der.} \\
\texttt{Giv nu endnu en kommasepareret liste med mindst 1 og op til 2 ord fra følgende ordliste, som bedst relaterer til eller er tættest forbundet med ordet 'ferie'. Husk at starte dit svar med GÆTS: . Bemærk: nogle ord er blevet fjernet fra listen sammenlignet med tidligere anmodninger.} \\
\\ 
\texttt{center, æble, model, stat, tjek, palme, pool, olympus, kirke, advokat, bold, blæksprutte, glas} \\
            }
        }
    }
    & \\ \\

    \theutterance \stepcounter{utterance}  
    & & & \multicolumn{4}{p{0.6\linewidth}}{
        \cellcolor[rgb]{0.9,0.9,0.9}{
            \makecell[{{p{\linewidth}}}]{
                \texttt{\tiny{[GM$\langle$P2]}}
                \texttt{GÆTS: palme, pool} \\
            }
        }
    }
     \\ \\

    \theutterance \stepcounter{utterance}  
    & & & \multicolumn{2}{p{0.3\linewidth}}{
        \cellcolor[rgb]{0.9,0.9,0.9}{
            \makecell[{{p{\linewidth}}}]{
                \texttt{\tiny{[GM$|$GM]}}
                ['palme', 'pool']
            }
        }
    }
    & & \\ \\

    \theutterance \stepcounter{utterance}  
    & & & \multicolumn{2}{p{0.3\linewidth}}{
        \cellcolor[rgb]{0.9,0.9,0.9}{
            \makecell[{{p{\linewidth}}}]{
                \texttt{\tiny{[GM$|$GM]}}
                {'word': 'palme', 'assignment': 'team'}
            }
        }
    }
    & & \\ \\

    \theutterance \stepcounter{utterance}  
    & & & \multicolumn{2}{p{0.3\linewidth}}{
        \cellcolor[rgb]{0.9,0.9,0.9}{
            \makecell[{{p{\linewidth}}}]{
                \texttt{\tiny{[GM$|$GM]}}
                {'word': 'palme', 'assignment': 'team'}
            }
        }
    }
    & & \\ \\

    \theutterance \stepcounter{utterance}  
    & & & \multicolumn{2}{p{0.3\linewidth}}{
        \cellcolor[rgb]{0.9,0.9,0.9}{
            \makecell[{{p{\linewidth}}}]{
                \texttt{\tiny{[GM$|$GM]}}
                {'word': 'pool', 'assignment': 'opponent'}
            }
        }
    }
    & & \\ \\

    \theutterance \stepcounter{utterance}  
    & & & \multicolumn{2}{p{0.3\linewidth}}{
        \cellcolor[rgb]{0.9,0.9,0.9}{
            \makecell[{{p{\linewidth}}}]{
                \texttt{\tiny{[GM$|$GM]}}
                \texttt{pool} \\
            }
        }
    }
    & & \\ \\

    \theutterance \stepcounter{utterance}  
    & & & \multicolumn{2}{p{0.3\linewidth}}{
        \cellcolor[rgb]{0.9,0.9,0.9}{
            \makecell[{{p{\linewidth}}}]{
                \texttt{\tiny{[GM$|$GM]}}
                {'word': 'olympus', 'assignment': 'opponent'}
            }
        }
    }
    & & \\ \\

    \theutterance \stepcounter{utterance}  
    & & & \multicolumn{2}{p{0.3\linewidth}}{
        \cellcolor[rgb]{0.9,0.9,0.9}{
            \makecell[{{p{\linewidth}}}]{
                \texttt{\tiny{[GM$|$GM]}}
                {'hidden': {'team': ['tjek'], 'innocent': ['model', 'æble', 'center'], 'opponent': ['glas'], 'assassin': ['advokat', 'kirke', 'blæksprutte', 'stat', 'bold']}, 'revealed': {'team': {'team': ['flagermus', 'and', 'krone', 'føniks', 'vakuum', 'jern', 'rør', 'palme'], 'innocent': [], 'opponent': ['presse', 'pool'], 'assassin': []}, 'opponent': {'team': [], 'innocent': [], 'opponent': ['server', 'sæl', 'hest', 'bjørn', 'olympus'], 'assassin': []}}}
            }
        }
    }
    & & \\ \\

    \theutterance \stepcounter{utterance}  
    & & \multicolumn{4}{p{0.6\linewidth}}{
        \cellcolor[rgb]{0.9,0.9,0.9}{
            \makecell[{{p{\linewidth}}}]{
                \texttt{\tiny{[P1$\langle$GM]}}
                \texttt{Ordet palme blev gættet korrekt. Ordet pool blev gættet, men det er et modstander ord. Dit holdkammerats tur endte der.} \\
\texttt{Giv nu endnu en ledetråd, der relaterer til nogle af dine resterende teamord, og en liste over de relaterede målord. Husk at starte din ledetråd med 'VEJLEDNING: ', sæt en ny linje, og start din kommaseparerede liste over målord med 'MÅL: '. Bemærk: Nogle ord er blevet fjernet fra listerne i forhold til tidligere anmodninger.} \\
\\ 
\texttt{Dine resterende holdord er: tjek.} \\
\texttt{Resterende ord for din modstander er: olympus, glas.} \\
\texttt{Resterende distraherende ord er: model, æble, center.} \\
\texttt{Resterende dræberord er: advokat, kirke, blæksprutte, stat, bold.} \\
            }
        }
    }
    & \\ \\

    \theutterance \stepcounter{utterance}  
    & \multicolumn{4}{p{0.6\linewidth}}{
        \cellcolor[rgb]{0.9,0.9,0.9}{
            \makecell[{{p{\linewidth}}}]{
                \texttt{\tiny{[P1$\rangle$GM]}}
                \texttt{VEJLEDNING: BE number.} \\
\\ 
\texttt{**P**} \\
\\ 
\texttt{{-} Run your response in the required format. The only remaining word is "tjek". The clue must relate to it.} \\
\texttt{{-} "tjek" means check (bank check), checkmark, or check (verification).} \\
\texttt{{-} Avoid: model, æble, center, advokat, kirke, blæksprutte, stat, bold.} \\
\texttt{{-} "underskrift" (signature) relates to check (bank) but maybe not checkmark.} \\
\texttt{{-} "hage" (checkmark) relates to checkmark but not bank check.} \\
\texttt{{-} "godkendelse" (approval) relates to verification.} \\
\texttt{{-} "betaling" (payment) relates to bank check.} \\
\texttt{{-} "kontrol" (control/check) relates to verification.} \\
\texttt{{-} "validering" (validation).} \\
\texttt{{-} "verificering".} \\
\texttt{{-} "bekræftelse" (confirmation).} \\
\texttt{{-} "godkendelse".} \\
\texttt{{-} "underskrift" is good for bank check.} \\
\texttt{{-} "hage" is the Danish word for the symbol ✓.} \\
\texttt{{-} "afkryssning".} \\
\texttt{{-} "liste" (list/checklist).} \\
\texttt{{-} "afhækning".} \\
\texttt{{-} "godkendt".} \\
\texttt{{-} "korrekt".} \\
\texttt{{-} "ok".} \\
\\ 
\texttt{"Godkendelse" covers verification.} \\
\texttt{"Underskrift" covers bank check.} \\
\texttt{"Hage" covers the symbol.} \\
\texttt{"Kontrol" covers the action.} \\
\\ 
\texttt{"Kontrol" is broad. "Godkendelse" is good. "Underskrift" is specific to banking. "Hage" is the symbol.} \\
\texttt{"Tjek" in Codenames context often clusters with "check" (bank), "check" (mark), "check" (verify).} \\
\texttt{"Kontrol" touches verify.} \\
\texttt{"Godkendelse" touches verify.} \\
\texttt{"Underskrift" touches bank check.} \\
\texttt{"Hage" touches mark.} \\
\\ 
\texttt{"Godkendelse" seems safest semantically for "tjek" (verify/check off).} \\
\texttt{"Underskrift" for check (bank).} \\
\texttt{"Kontrol" for check (inspect).} \\
\\ 
\texttt{I'll go with "GODKENDELSE" as it implies a checkmark or verification. Or "UNDERSKRIFT" for cheque.} \\
\texttt{"Tjek" is often "check" (English loan) or "check" (mark).} \\
\texttt{"Godkendelse" {-}> Tjek (som} \\
            }
        }
    }
    & & \\ \\

    \theutterance \stepcounter{utterance}  
    & & & \multicolumn{2}{p{0.3\linewidth}}{
        \cellcolor[rgb]{0.9,0.9,0.9}{
            \makecell[{{p{\linewidth}}}]{
                \texttt{\tiny{[GM$|$GM]}}
                {'player': 'Cluegiver', 'type': 'rambling error', 'utterance': 'VEJLEDNING: BE number.\n\n**P**\n\n- Run your response in the required format. The only remaining word is "tjek". The clue must relate to it.\n- "tjek" means check (bank check), checkmark, or check (verification).\n- Avoid: model, æble, center, advokat, kirke, blæksprutte, stat, bold.\n- "underskrift" (signature) relates to check (bank) but maybe not checkmark.\n- "hage" (checkmark) relates to checkmark but not bank check.\n- "godkendelse" (approval) relates to verification.\n- "betaling" (payment) relates to bank check.\n- "kontrol" (control/check) relates to verification.\n- "validering" (validation).\n- "verificering".\n- "bekræftelse" (confirmation).\n- "godkendelse".\n- "underskrift" is good for bank check.\n- "hage" is the Danish word for the symbol ✓.\n- "afkryssning".\n- "liste" (list/checklist).\n- "afhækning".\n- "godkendt".\n- "korrekt".\n- "ok".\n\n"Godkendelse" covers verification.\n"Underskrift" covers bank check.\n"Hage" covers the symbol.\n"Kontrol" covers the action.\n\n"Kontrol" is broad. "Godkendelse" is good. "Underskrift" is specific to banking. "Hage" is the symbol.\n"Tjek" in Codenames context often clusters with "check" (bank), "check" (mark), "check" (verify).\n"Kontrol" touches verify.\n"Godkendelse" touches verify.\n"Underskrift" touches bank check.\n"Hage" touches mark.\n\n"Godkendelse" seems safest semantically for "tjek" (verify/check off).\n"Underskrift" for check (bank).\n"Kontrol" for check (inspect).\n\nI\'ll go with "GODKENDELSE" as it implies a checkmark or verification. Or "UNDERSKRIFT" for cheque.\n"Tjek" is often "check" (English loan) or "check" (mark).\n"Godkendelse" -> Tjek (som', 'message': 'Your answer contained more than two lines, please only give one clue and your targets on two separate lines.'}
            }
        }
    }
    & & \\ \\

\end{supertabular}
}

\end{document}
